\documentclass[14pt]{article}
\usepackage[top=25truemm,bottom=20truemm,left=25truemm,right=25truemm]{geometry}
\usepackage{xskak}
\usepackage{fancyhdr}
\setlength\parindent{0pt}

\begin{document}

\begin{titlepage}
\thispagestyle{fancy}

\lhead{\Large ADMIN ONLY \hspace{3mm} Time submitted: \underline{\hspace{2.5cm}} Total score: \underline{\hspace{2.5cm}}/100}

\begin{center}
\vspace{20mm}
\textbf{\Huge 2026 AUSTRALIAN JUNIOR}\\
\vspace{6mm}
\textbf{\Huge CHESS CHAMPIONSHIP}\\
\vspace{6mm}
\textbf{\Huge PROBLEM SOLVING ALPHA}\\
\vspace{6mm}
\textbf{\Huge CANBERRA}\\
\end{center}

\vspace{8mm}\\
\textbf{\LARGE NAME: \hspace{3mm} \underline{\hspace{8cm}}}\\
\vspace{5mm}\\
\textbf{\LARGE DATE OF BIRTH: \hspace{3mm} \underline{\hspace{5cm}}}\\
\vspace{5mm}\\

\textbf{\LARGE Please circle your age group:}\\
\begin{itemize}
  \item \textbf{\Large  U8 (born in 2019 or after)    U14 (born in 2013 or after)}
  \item \textbf{\Large  U10 (born in 2017 or after)   U16 (born in 2011 or after)}
  \item \textbf{\Large  U12 (born in 2015 or after)   U18 (born in 2009 or after)}
\end{itemize}
\vspace{5mm}\\

\textbf{\LARGE Please circle: Boy or Girl}\\
\vspace{5mm}\\

\textbf{\LARGE Please read the following rules.}\\
\begin{itemize}
  \item {\Large You have up to 2 hours for the problem solving paper. You may not need all that time. You may leave after handing in your answers.}
  \item {\Large The questions generally go from easy to hard, so it’s best to start from the beginning and go in order.}
  \item {\Large You can use your chess board and pieces to help you.}
  \item {\Large Each page show a question and the marks you can score. Please write your answers in the space under each diagram. Altogether, the test has 100 marks. Even if you don’t give perfect answers, you can still get some partial marks.}
  \item {\Large For endgame puzzles, show enough moves to reach a position where you are clearly winning or drawing.}
  \item {\Large You don’t have to finish them all, do your best and solve as many puzzles as you can.}
\end{itemize}
\vspace{3mm}\\
\textbf{\LARGE Enjoy the problem solving.}
\end{titlepage}


\newpage
\textbf{\Large Example Q1. Mate in 2.}\\

{\setchessboard{boardfontsize=20pt,labelfontsize=12pt}
\newchessgame
\chessboard[setfen=k7/7R/2K5/8/8/8/8/8 w - - 0 1]}
\hspace{3mm} \Large{The open box shows White's move.}\\

\textbf{\Large Answer:}\\

\Large{1. Kb6  \hspace{5mm} Kb8 \hspace{8mm} \textleftarrow \hspace{2mm} Write the first White move, then the first Black move.}\\
\Large{2. Rh8\# (1-0) \hspace{8mm} \textleftarrow \hspace{2mm} Write the second White move, which is a mate.}\\

\textbf{\Large Example Q2. Black moves and draws.}\\
{\setchessboard{inverse, boardfontsize=20pt,labelfontsize=12pt}
\newchessgame
\chessboard[setfen=3k4/8/3PK3/8/8/8/8/8 b - - 0 1]}
\hspace{3mm} \Large{The closed box shows Black's move.}\\

\textbf{\Large Answer:}\\

\Large{1. ... \hspace{5mm} Ke8}\\
\Large{2. d7+ \hspace{2mm} Kd8}\\
\Large{3. Kd6 (draw)  \hspace{8mm} \textleftarrow \hspace{2mm} Write enough moves showing a clear draw position.}\\


\newpage
\textbf{\Large Q1. Mate in 2. (6 marks)}\\
{\setchessboard{boardfontsize=30pt,labelfontsize=14pt}
\newchessgame
\chessboard[setfen=8/7R/k1Kp4/3N4/r7/2Q5/8/8 w - - 0 1]}\\

\textbf{\Large Answer:}\\


\vfill{\Large Source: Jeff Coakley, Winning Chess Puzzles for Kids Vol. 2}


\newpage
\textbf{\Large Q2. Mate in 2 (6 marks)}\\
{\setchessboard{boardfontsize=30pt,labelfontsize=14pt}
\newchessgame
\chessboard[setfen=8/2p5/4K3/8/3pkp2/3ppp2/8/4Q3 w - - 0 1]}\\

\textbf{\Large Answer:}\\


\vfill{\Large Source: Thomas Salthouse, Falkirk Herald, 1914 (YACPDB: 22264)}


\newpage
\textbf{\Large Q3. Mate in 2. (6 marks)}\\
{\setchessboard{boardfontsize=30pt,labelfontsize=14pt}
\newchessgame
\chessboard[setfen=8/p1K5/krP2Q2/NRR5/8/8/8/8 w - - 0 1]}\\

\textbf{\Large Answer:}\\


\vfill\Large {Source: Cornelis Goldschmeding, problem (Zagreb), 1957 (YACPDB: 8721)}


\newpage
\textbf{\Large Q4. Mate in 2. (6 marks)}\\
{\setchessboard{boardfontsize=30pt,labelfontsize=14pt}
\newchessgame
\chessboard[setfen=kbK5/pp6/1P6/8/8/8/8/R7 w - - 0 1]}\\

\textbf{\Large Answer:}\\


\vfill{\Large Source: Paul Morphy, 1910 (YACPDB: 625426)}


\newpage
\textbf{\Large Q5. Mate in 2. (6 marks)}\\
{\setchessboard{boardfontsize=30pt,labelfontsize=14pt}
\newchessgame
\chessboard[setfen=8/8/8/8/4R3/6k1/8/4K2R w K - 0 1]}\\

\textbf{\Large Answer:}\\


\vfill{\Large Source: W. E. Candy, 1911 (YACPDB: 3889)}


\newpage
\textbf{\Large Q6. What is the missing piece? Whose turn is it? Which square does the missing piece belong on? Show the previous moves that reached this position. (8 marks)}\\
{\setchessboard{boardfontsize=30pt, labelfontsize=14pt, showmover=false}
\newchessgame
\chessboard[setfen=8/8/8/1r1b4/B7/8/8/3k4 - - 0 1]}\\

\textbf{\Large Answer:}\\

\vfill{\Large Source: Raymond M. Smullyan: The Chess Mysteries of the Arabian Knights}


\newpage
\textbf{\Large Q7. WHO moved last? WHICH piece was moved from WHERE to WHERE? (6 marks)}\\
{\setchessboard{boardfontsize=30pt, labelfontsize=14pt, showmover=false}
\newchessgame
\chessboard[setfen=rnR1kbnr/pp2pppp/3p4/1B4N1/1R1BP3/3P2P1/PPP2PP1/3Q1KN1 b kq - 0 1]}\\

\textbf{\Large Answer:}\\


\vfill{\Large Source: Noam Elkies, U.S. Problem Bulletin, 1995}


\newpage
\textbf{\Large Q8. The position below occured after the Black's 4th move, from the starting position. Show the moves to reach this position. (6 marks)}\\
{\setchessboard{boardfontsize=30pt,labelfontsize=14pt}
\newchessgame
\chessboard[setfen=rnbqkbnr/pp3ppp/2p1p3/8/4P3/8/PPPP1PPP/RNBQK1NR]}\\

\textbf{\Large Answer:}\\


\vfill{\Large Source: Tibor Orbán, Die Schwalbe, 1976}


\newpage
\textbf{\Large Q9. The position below occured after the Black's 4th move, from the starting position. Show the moves to reach this position. (6 marks)}\\
{\setchessboard{boardfontsize=30pt,labelfontsize=14pt}
\newchessgame
\chessboard[setfen=rnbqkb1r/pppppp2/6p1/7p/3P4/8/PPP1PPPP/RN2KBNR w KQkq - 0 5]}\\

\textbf{\Large Answer:}\\


\vfill{\Large Source: L. Mano, source unknown, 1999}


\newpage
\textbf{\Large Q10. Helpmate in 1. What is Black's move to allow White a  "Mate in 1"? Answer Black move's first, then White's checkmates. (6 marks)}\\
{\setchessboard{boardfontsize=30pt,labelfontsize=14pt}
\newchessgame
\chessboard[setfen=r3k1nr/pp2q2p/1n2pb1B/1N1b1p2/8/P5N1/1P2QPPP/3RR1K1 b kq - 0 1]}\\

\textbf{\Large Answer:}\\


\vfill{\Large Source: Jeff Coakley, Winning Chess Puzzles for Kids Vol. 2}


\newpage
\textbf{\Large 11. Helpmate in 1. Black's positions are all different, but White's positions are all same as the previous question. (6 marks each)}\\
\textbf{\Large Q11A.}
{\setchessboard{boardfontsize=20pt,labelfontsize=12pt}
\newchessgame
\chessboard[setfen=r2qk2r/pp1n3p/4pb1B/1N1b1p2/8/P5N1/1P2QPPP/3RR1K1 b kq - 0 1]}
\textbf{\Large Answer:}\\
\textbf{\Large Q11B.}
{\setchessboard{boardfontsize=20pt,labelfontsize=12pt}
\newchessgame
\chessboard[setfen=r2qk2r/pp2n2p/1n2pb1B/1N1b1p2/8/P5N1/1P2QPPP/3RR1K1 b kq - 1 2]}
\textbf{\Large Answer:}\\
\textbf{\Large Q11C.}
{\setchessboard{boardfontsize=20pt,labelfontsize=12pt}
\newchessgame
\chessboard[setfen=r2qk1nr/pp5p/1n2pb1B/1N1b1p2/8/P5N1/1P2QPPP/3RR1K1 b kq - 2 3]}
\textbf{\Large Answer:}\\
\vfill{\Large Source: Jeff Coakley, Winning Chess Puzzles for Kids Vol. 2}


\newpage
\textbf{\Large Q12. Helpmate in 3. Black moves 1st and White checkmates with their 3rd move. (6 marks)}\\
{\setchessboard{boardfontsize=30pt,labelfontsize=14pt}
\newchessgame
\chessboard[setfen=1K6/8/1N6/1k6/8/1B6/8/8 b - - 0 1]}\\

\textbf{\Large Answer:}\\


\vfill{\Large Source: Hilmar Ebert, Schach-Echo, 1977 (YACPDB: 638060)}


\newpage
\textbf{\Large Q13. Selfmate in 2. White forces Black to checkmate White's King with Black's 2nd move. (NB This must be a forced variation as Black will try and avoid checkmating White if possible!) (6 marks)}\\
{\setchessboard{boardfontsize=30pt,labelfontsize=14pt}
\newchessgame
\chessboard[setfen=K1k1r3/P7/7Q/8/8/8/8/8 w - - 0 1]}\\

\textbf{\Large Answer:}\\


\vfill{\Large Source: Wilhelm Krämer, Die Schwalbe, 1928 (YACPDB: 106430)}


\newpage
\textbf{\Large Q14. Selfmate in 2. White forces Black to checkmate White in 2 moves. (8 marks)}\\
{\setchessboard{boardfontsize=30pt,labelfontsize=14pt}
\newchessgame
\chessboard[setfen=KB3N2/P1P1p1P1/5P1k/4P2p/7P/8/6B1/7b w - - 0 1]}\\

\textbf{\Large Answer:}\\


\vfill{\Large Source: Wolfgang Pauly, The Theory of Pawn Promotion, 1912 (YACPDB: 47304)}


\newpage
\textbf{\Large Q15. White to play. Who wins?}\\
{\setchessboard{boardfontsize=30pt,labelfontsize=14pt}
\newchessgame
\chessboard}\\

\textbf{\Large Answer:}\\


\vfill{\Large Source: Composer unknown, date unknown}

\end{document}
